\paragraph{位置--尺度 Poisson--Cauchy 粗粒化的正规形、熵层级与端点重写}\label{par:xi-poisson-cauchy-location-scale-entropy-rigidity-busemann}
在定理 \ref{thm:xi-cauchy-poisson-two-channel-orthogonal-leading} 与定理 \ref{thm:xi-poisson-cauchy-harmonic-moment-multipole-kl} 的框架内，
位置--尺度联合混合可被统一压缩为“二阶协方差不变量 \(\to\) 熵主系数 \(\to\) 端点几何坐标”的闭合链。
下述命题组给出该闭合链的显式形式，并把二阶盲点与六次衰减写为可核验的刚性判别式。

\begin{proposition}[位置--尺度混合的二阶正规形与调和塌缩]\label{prop:xi-poisson-cauchy-location-scale-second-normal-form}
\leanverified{paper\_xi\_poisson\_cauchy\_location\_scale\_second\_normal\_form}
令
\[
\mathsf P_t(x):=\frac{1}{\pi}\frac{t}{x^2+t^2}\qquad(t>0),
\]
并设 \(\tilde\nu\) 为 \(\RR\times(0,\infty)\) 上概率测度。
定义
\[
\tilde h_t(x):=\int_{\RR\times(0,\infty)} \mathsf P_{t+\delta}(x-\gamma)\,\dd\tilde\nu(\gamma,\delta).
\]
记
\[
\bar\gamma:=\int\gamma\,\dd\tilde\nu,\qquad
\bar\delta:=\int\delta\,\dd\tilde\nu,\qquad
\tilde g_t(x):=\mathsf P_{t+\bar\delta}(x-\bar\gamma),
\]
并定义二阶不变量
\[
A:=\Var_{\tilde\nu}(\gamma)-\Var_{\tilde\nu}(\delta),\qquad
B:=\Cov_{\tilde\nu}(\gamma,\delta).
\]
若三阶中心绝对矩有限，则存在余项 \(\varepsilon_t\) 使
\[
\tilde h_t
=\tilde g_t+\frac{A}{2}\,\partial_x^2\tilde g_t-B\,\partial_x\partial_t\tilde g_t+\varepsilon_t,
\qquad
\|\varepsilon_t\|_{L^1(\RR)}=o\!\bigl((t+\bar\delta)^{-2}\bigr).
\]
\end{proposition}

\begin{proof}
令 \(s=t+\bar\delta\)、\(z=x-\bar\gamma\)，并在 \(\tilde\nu\) 下置
\[
U:=\gamma-\bar\gamma,\qquad V:=\delta-\bar\delta.
\]
则
\[
\tilde h_t(x)=\mathbb E\!\left[\mathsf P_{s+V}(z-U)\right],\qquad
\tilde g_t(x)=\mathsf P_s(z).
\]
对 \((u,v)\mapsto \mathsf P_{s+v}(z-u)\) 在 \((0,0)\) 作二阶 Taylor 展开并取期望，线性项因 \(\mathbb E U=\mathbb E V=0\) 消失，得到
\[
\tilde h_t-\tilde g_t
=\frac{\mathbb E[U^2]}{2}\,\partial_z^2\mathsf P_s
-\mathbb E[UV]\,\partial_z\partial_s\mathsf P_s
+\frac{\mathbb E[V^2]}{2}\,\partial_s^2\mathsf P_s+\varepsilon_t.
\]
Poisson 核满足调和关系
\[
\partial_s^2\mathsf P_s+\partial_z^2\mathsf P_s=0,
\]
故 \(\partial_s^2\mathsf P_s=-\partial_z^2\mathsf P_s\)，主项合并为
\[
\frac{\Var(U)-\Var(V)}{2}\,\partial_z^2\mathsf P_s-\Cov(U,V)\,\partial_z\partial_s\mathsf P_s
=\frac{A}{2}\,\partial_x^2\tilde g_t-B\,\partial_x\partial_t\tilde g_t.
\]
三阶余项由三阶中心绝对矩与 \(\mathsf P_s\) 三阶导数的 \(L^1\)-尺度估计控制，故
\(\|\varepsilon_t\|_{L^1}=o(s^{-2})\)。证毕。
\end{proof}

\begin{theorem}[二阶熵系数的锐极限与复二阶矩分解]\label{thm:xi-poisson-cauchy-location-scale-kl-sharp-m2}
在命题 \ref{prop:xi-poisson-cauchy-location-scale-second-normal-form} 条件下，
若四阶中心矩有限，则有
\[
\lim_{t\to\infty}(t+\bar\delta)^4 D_{\mathrm{KL}}(\tilde h_t\mid \tilde g_t)
=\frac{A^2}{8}+\frac{B^2}{2}.
\]
若定义
\[
Z:=(\gamma-\bar\gamma)+i(\delta-\bar\delta),
\]
则等价地
\[
\lim_{t\to\infty}(t+\bar\delta)^4 D_{\mathrm{KL}}(\tilde h_t\mid \tilde g_t)
=\frac18\,\bigl|\mathbb E_{\tilde\nu}[Z^2]\bigr|^2,
\qquad
\mathbb E_{\tilde\nu}[Z^2]=A+2iB.
\]
\end{theorem}

\begin{proof}
令 \(s=t+\bar\delta\)、\(y=(x-\bar\gamma)/s\)、\(G(y):=\pi^{-1}(1+y^2)^{-1}\)。
由命题 \ref{prop:xi-poisson-cauchy-location-scale-second-normal-form}，
\[
\frac{\tilde h_t}{\tilde g_t}=1+\delta_t,\qquad
\delta_t(x)=\frac{1}{s^2}\bigl(Au(y)-Bv(y)\bigr)+o_{L^2(\tilde g_t\dd x)}(s^{-2}),
\]
其中
\[
u(y)=\frac{3y^2-1}{(1+y^2)^2}\ \text{(偶)},\qquad
v(y)=\frac{2y(3-y^2)}{(1+y^2)^2}\ \text{(奇)}.
\]
由 \(\int \tilde g_t\delta_t\,\dd x=0\) 与
\((1+\delta)\log(1+\delta)=\delta+\frac12\delta^2+O(\delta^3)\)，
\[
D_{\mathrm{KL}}(\tilde h_t\mid \tilde g_t)
=\frac12\int \tilde g_t\,\delta_t^2\,\dd x+o(s^{-4}).
\]
代入主项并利用奇偶正交性 \(\int_{\RR}Guv\,\dd y=0\)，得
\[
D_{\mathrm{KL}}(\tilde h_t\mid \tilde g_t)
=\frac{1}{2s^4}\!\left(
A^2\!\int_{\RR}Gu^2\,\dd y+B^2\!\int_{\RR}Gv^2\,\dd y\right)+o(s^{-4}).
\]
标准积分给出
\[
\int_{\RR}Gu^2\,\dd y=\frac14,\qquad
\int_{\RR}Gv^2\,\dd y=1,
\]
故结论成立。再由
\[
\mathbb E_{\tilde\nu}[Z^2]
=\mathbb E\!\left[(\gamma-\bar\gamma)^2-(\delta-\bar\delta)^2\right]
+2i\,\mathbb E\!\left[(\gamma-\bar\gamma)(\delta-\bar\delta)\right]
=A+2iB
\]
得复矩形式。证毕。
\end{proof}

\begin{theorem}[Cauchy 位置--尺度族的信息投影与二分常数律]\label{thm:xi-poisson-cauchy-location-scale-information-projection}
\leanverified{paper\_xi\_poisson\_cauchy\_location\_scale\_information\_projection}
沿用命题 \ref{prop:xi-poisson-cauchy-location-scale-second-normal-form} 的记号。
对任意 \((\alpha,\beta)\in\RR^2\)，定义一族位置--尺度可调参考密度
\begin{equation}\label{eq:xi-poisson-cauchy-location-scale-gab}
\tilde g_t^{(\alpha,\beta)}(x)
:=\mathsf P_{\,t+\bar\delta+\alpha/(t+\bar\delta)}\!\left(x-\bar\gamma-\beta/(t+\bar\delta)\right).
\end{equation}
则存在唯一一对 \((\alpha^\star,\beta^\star)\) 使得
\begin{equation}\label{eq:xi-poisson-cauchy-location-scale-information-projection-optimal}
\lim_{t\to\infty}(t+\bar\delta)^4\,D_{\mathrm{KL}}\!\left(\tilde h_t\mid \tilde g_t^{(\alpha^\star,\beta^\star)}\right)
=\inf_{\alpha,\beta}\ \limsup_{t\to\infty}(t+\bar\delta)^4\,D_{\mathrm{KL}}\!\left(\tilde h_t\mid \tilde g_t^{(\alpha,\beta)}\right),
\end{equation}
并且最优参数与最小极限值满足闭式
\[
\alpha^\star=\frac{A}{2},\qquad
\beta^\star=-B,
\]
\[
\lim_{t\to\infty}(t+\bar\delta)^4\,D_{\mathrm{KL}}\!\left(\tilde h_t\mid \tilde g_t^{(\alpha^\star,\beta^\star)}\right)
=\frac{1}{16}\Bigl(A^2+4B^2\Bigr)
=\frac{A^2}{16}+\frac{B^2}{4}.
\]
\end{theorem}
\begin{proof}
记 \(s:=t+\bar\delta\)、\(y:=(x-\bar\gamma)/s\)、\(G(y):=\pi^{-1}(1+y^2)^{-1}\)。
由定理 \ref{thm:xi-poisson-cauchy-location-scale-kl-sharp-m2} 证明中的展开，有
\[
\frac{\tilde h_t}{\tilde g_t}=1+\frac{1}{s^2}\bigl(Au(y)-Bv(y)\bigr)+o_{L^2(\tilde g_t\dd x)}(s^{-2}),
\]
其中
\[
u(y)=\frac{3y^2-1}{(1+y^2)^2}\ \text{(偶)},\qquad
v(y)=\frac{2y(3-y^2)}{(1+y^2)^2}\ \text{(奇)}.
\]
\par
另一方面，对 \eqref{eq:xi-poisson-cauchy-location-scale-gab} 作二阶尺度的相对扰动展开。
记
\[
\phi(y):=\frac{y^2-1}{1+y^2}\ \text{(偶)},\qquad
\psi(y):=\frac{2y}{1+y^2}\ \text{(奇)}.
\]
则直接由位置与尺度方向的得分函数得到
\[
\frac{\tilde g_t^{(\alpha,\beta)}}{\tilde g_t}
=1+\frac{1}{s^2}\bigl(\alpha\,\phi(y)+\beta\,\psi(y)\bigr)+o_{L^2(\tilde g_t\dd x)}(s^{-2}),
\]
从而
\[
\frac{\tilde h_t}{\tilde g_t^{(\alpha,\beta)}}-1
=\frac{1}{s^2}\Bigl(Au-\alpha\phi-Bv-\beta\psi\Bigr)+o_{L^2(\tilde g_t\dd x)}(s^{-2}).
\]
由 \(\int \tilde g_t\left(\frac{\tilde h_t}{\tilde g_t^{(\alpha,\beta)}}-1\right)\dd x=0\) 与 \((1+z)\log(1+z)=z+\frac12z^2+O(z^3)\)，得
\[
D_{\mathrm{KL}}(\tilde h_t\mid \tilde g_t^{(\alpha,\beta)})
=\frac{1}{2s^4}\int_{\RR}G(y)\Bigl(Au-\alpha\phi-Bv-\beta\psi\Bigr)^2\dd y+o(s^{-4}).
\]
利用偶奇正交性 \(\int Gu\psi=\int Gv\phi=0\)，该二次型分裂为
\[
\int G(Au-\alpha\phi-Bv-\beta\psi)^2
=\int G(Au-\alpha\phi)^2+\int G(Bv+\beta\psi)^2.
\]
令内积 \(\langle f,g\rangle:=\int_{\RR}Gfg\)。
直接计算得
\[
\langle u,u\rangle=\frac14,\qquad
\langle v,v\rangle=1,\qquad
\langle u,\phi\rangle=\frac14,\qquad
\langle \phi,\phi\rangle=\frac12,\qquad
\langle v,\psi\rangle=\frac12,\qquad
\langle \psi,\psi\rangle=\frac12.
\]
于是两块分别在
\[
\alpha^\star=A\frac{\langle u,\phi\rangle}{\langle \phi,\phi\rangle}=\frac{A}{2},
\qquad
\beta^\star=-B\frac{\langle v,\psi\rangle}{\langle \psi,\psi\rangle}=-B
\]
处唯一取极小；对应的最小值为
\[
\inf_\alpha \int G(Au-\alpha\phi)^2
=A^2\left(\langle u,u\rangle-\frac{\langle u,\phi\rangle^2}{\langle \phi,\phi\rangle}\right)
=\frac{A^2}{8},
\]
\[
\inf_\beta \int G(Bv+\beta\psi)^2
=B^2\left(\langle v,v\rangle-\frac{\langle v,\psi\rangle^2}{\langle \psi,\psi\rangle}\right)
=\frac{B^2}{2}.
\]
代回得到
\[
D_{\mathrm{KL}}(\tilde h_t\mid \tilde g_t^{(\alpha^\star,\beta^\star)})
=\frac{1}{2s^4}\left(\frac{A^2}{8}+\frac{B^2}{2}\right)+o(s^{-4})
=\frac{A^2+4B^2}{16s^4}+o(s^{-4}),
\]
从而结论成立。唯一性由二次型严格凸性给出。证毕。
\end{proof}

\begin{corollary}[一般光滑 \texorpdfstring{$f$}{f}-散度的同伦二次型与投影常数]\label{cor:xi-poisson-cauchy-location-scale-fdiv-projection-constants}
在命题 \ref{prop:xi-poisson-cauchy-location-scale-second-normal-form} 的设定下，设 \(f:(0,\infty)\to\RR\) 在 \(1\) 的邻域 \(C^2\)，满足
\[
f(1)=f'(1)=0,\qquad f''(1)>0,
\]
并定义 \(f\)-散度
\[
D_f(h\mid g):=\int_{\RR} g(x)\,f\!\left(\frac{h(x)}{g(x)}\right)\dd x.
\]
则：
\begin{enumerate}
  \item 若取固定一阶匹配基线 \(g=\tilde g_t\)，则
  \[
  \lim_{t\to\infty}(t+\bar\delta)^4\,D_f(\tilde h_t\mid \tilde g_t)
  =\frac{f''(1)}{8}\Bigl(A^2+4B^2\Bigr).
  \]
  \item 若取位置--尺度投影基线 \(g=\tilde g_t^{(\alpha^\star,\beta^\star)}\)（定理 \ref{thm:xi-poisson-cauchy-location-scale-information-projection}），则
  \[
  \lim_{t\to\infty}(t+\bar\delta)^4\,D_f\!\left(\tilde h_t\mid \tilde g_t^{(\alpha^\star,\beta^\star)}\right)
  =\frac{f''(1)}{16}\Bigl(A^2+4B^2\Bigr).
  \]
\end{enumerate}
\end{corollary}
\begin{proof}
由 \(f\) 在 \(1\) 处的二阶展开
\(
f(1+z)=\frac{f''(1)}{2}z^2+o(z^2)
\)（\(z\to 0\)），并用
\(
\delta_t:=\tilde h_t/\tilde g_t-1=O_{L^2(\tilde g_t\dd x)}((t+\bar\delta)^{-2})
\)
以及 \(\int \tilde g_t\delta_t\,\dd x=0\)，可得
\[
D_f(\tilde h_t\mid \tilde g_t)=\frac{f''(1)}{2}\int \tilde g_t\,\delta_t^2\,\dd x+o((t+\bar\delta)^{-4}).
\]
再由定理 \ref{thm:xi-poisson-cauchy-location-scale-kl-sharp-m2} 的主项常数
\(
\lim (t+\bar\delta)^4D_{\mathrm{KL}}(\tilde h_t\mid \tilde g_t)=\frac18(A^2+4B^2)
\)
与 \(D_{\mathrm{KL}}=\frac12\int \tilde g_t\delta_t^2+o((t+\bar\delta)^{-4})\)，得到第一条。
第二条同理由定理 \ref{thm:xi-poisson-cauchy-location-scale-information-projection} 的最小主项常数给出。证毕。
\end{proof}

\begin{corollary}[二阶盲点后的六次衰减常数]\label{cor:xi-poisson-cauchy-location-scale-m2-blind-m3-kl}
\leanverified{paper\_xi\_poisson\_cauchy\_location\_scale\_m2\_blind\_m3\_kl}
在定理 \ref{thm:xi-poisson-cauchy-location-scale-kl-sharp-m2} 设定下，进一步假设六阶中心矩有限。
若 \(\mathbb E_{\tilde\nu}[Z^2]=0\)，则
\[
\lim_{t\to\infty}(t+\bar\delta)^6 D_{\mathrm{KL}}(\tilde h_t\mid \tilde g_t)
=\frac{3}{32}\,\bigl|\mathbb E_{\tilde\nu}[Z^3]\bigr|^2.
\]
\end{corollary}

\begin{proof}
由 \(\mathbb E_{\tilde\nu}[Z^2]=0\) 得二阶多极消失。
应用定理 \ref{thm:xi-poisson-cauchy-harmonic-moment-multipole-kl} 并取 \(k=3\)（其常数 \(c_3=\binom42/2^6=3/32\)），即得结论。
证毕。
\end{proof}

\begin{corollary}[超四次与超六次衰减的层级刚性]\label{cor:xi-poisson-cauchy-location-scale-hierarchy-rigidity}
在上述条件下有：
\[
D_{\mathrm{KL}}(\tilde h_t\mid \tilde g_t)=o\!\bigl((t+\bar\delta)^{-4}\bigr)
\Longrightarrow \mathbb E_{\tilde\nu}[Z^2]=0,
\]
且在 \(\mathbb E_{\tilde\nu}[Z^2]=0\) 条件下，
\[
D_{\mathrm{KL}}(\tilde h_t\mid \tilde g_t)=o\!\bigl((t+\bar\delta)^{-6}\bigr)
\Longrightarrow \mathbb E_{\tilde\nu}[Z^3]=0.
\]
\end{corollary}

\begin{proof}
第一条由定理 \ref{thm:xi-poisson-cauchy-location-scale-kl-sharp-m2} 的极限常数立即推出。
第二条由推论 \ref{cor:xi-poisson-cauchy-location-scale-m2-blind-m3-kl} 的极限常数立即推出。
证毕。
\end{proof}

\begin{corollary}[耗散量的两级锐渐近常数]\label{cor:xi-poisson-cauchy-location-scale-i1-two-level}
在命题 \ref{prop:xi-poisson-kl-dissipation-fractional-fisher} 的耗散恒等式
\[
\frac{\dd}{\dd t}D_{\mathrm{KL}}(\tilde h_t\mid \tilde g_t)=-I_1(\tilde h_t\mid \tilde g_t)
\]
下，有
\[
\lim_{t\to\infty}(t+\bar\delta)^5 I_1(\tilde h_t\mid \tilde g_t)
=\frac12\,\bigl|\mathbb E_{\tilde\nu}[Z^2]\bigr|^2.
\]
若再满足 \(\mathbb E_{\tilde\nu}[Z^2]=0\)，则
\[
\lim_{t\to\infty}(t+\bar\delta)^7 I_1(\tilde h_t\mid \tilde g_t)
=\frac{9}{16}\,\bigl|\mathbb E_{\tilde\nu}[Z^3]\bigr|^2.
\]
\end{corollary}

\begin{proof}
由定理 \ref{thm:xi-poisson-cauchy-location-scale-kl-sharp-m2} 记
\(
D_{\mathrm{KL}}=C_4 (t+\bar\delta)^{-4}+o((t+\bar\delta)^{-4})
\)
且
\(
C_4=\frac18|\mathbb E_{\tilde\nu}[Z^2]|^2
\)；
对 \(t\) 求导即得
\(
I_1\sim 4C_4 (t+\bar\delta)^{-5}
\)，故首式成立。
同理，推论 \ref{cor:xi-poisson-cauchy-location-scale-m2-blind-m3-kl} 给出
\(
D_{\mathrm{KL}}=C_6 (t+\bar\delta)^{-6}+o((t+\bar\delta)^{-6})
\)
且
\(
C_6=\frac{3}{32}|\mathbb E_{\tilde\nu}[Z^3]|^2
\)，
故
\(
I_1\sim 6C_6 (t+\bar\delta)^{-7}
\)
并得第二式。证毕。
\end{proof}

\begin{corollary}[端点 Poisson--Busemann 坐标下的二阶熵系数重写]\label{cor:xi-poisson-cauchy-location-scale-busemann-rewrite}
令 \(w=\mathcal C(\gamma+i\delta)\in\mathbb D\)，其中 \(\mathcal C\) 为 Cayley 紧化（\eqref{eq:app-cayley-disk}）。
在 \(w\)-推前测度下定义
\[
X(w):=\Re\!\bigl(\mathcal C^{-1}(w)\bigr),\qquad
Y(w):=P_w(-1).
\]
则
\[
\lim_{t\to\infty}(t+\bar\delta)^4 D_{\mathrm{KL}}(\tilde h_t\mid \tilde g_t)
=\frac18\bigl(\Var(X)-\Var(Y)\bigr)^2+\frac12\,\Cov(X,Y)^2.
\]
\end{corollary}

\begin{proof}
由结论 \ref{con:xi-timefiber-poisson-busemann-identity}，
\(
Y(w)=P_w(-1)=\delta
\)；
而
\(
X(w)=\Re(\mathcal C^{-1}(w))=\gamma
\)。
故定理 \ref{thm:xi-poisson-cauchy-location-scale-kl-sharp-m2} 中的
\(
A=\Var(\gamma)-\Var(\delta)
\) 与
\(
B=\Cov(\gamma,\delta)
\)
可直接重写为
\(
\Var(X)-\Var(Y)
\) 与
\(
\Cov(X,Y)
\)。
代回即得结论。证毕。
\end{proof}

\begin{theorem}[Cauchy 位置--尺度族中的二阶最优尺度重整化]\label{thm:xi-cauchy-poisson-fixed-location-scale-optimal-renormalization}
沿用命题 \ref{prop:xi-poisson-tv-kl-second-order} 的单通道设定：给定 \(\tilde\nu\in\mathcal P(\RR)\)，令
\[
h_t(x):=\int_{\RR}\frac1\pi\frac{t}{(x-\gamma)^2+t^2}\,\dd\tilde\nu(\gamma),\qquad
\bar\gamma:=\int_{\RR}\gamma\,\dd\tilde\nu(\gamma),
\]
\[
\sigma^2:=\int_{\RR}(\gamma-\bar\gamma)^2\,\dd\tilde\nu(\gamma)\in(0,\infty).
\]
并假设四阶中心绝对矩有限：
\[
\int_{\RR}|\gamma-\bar\gamma|^4\,\dd\tilde\nu(\gamma)<\infty.
\]
对 \(\tau>0\) 定义固定位置的 Cauchy 族
\[
\mathcal C_{\bar\gamma,\tau}(x):=\frac1\pi\frac{\tau}{(x-\bar\gamma)^2+\tau^2},
\qquad
g_{t,\tau}(x):=\mathcal C_{\bar\gamma,\tau}(x),
\]
并记
\[
\mathcal D_f(t,\tau):=D_f(h_t\mid g_{t,\tau}),
\]
其中 \(f\) 凸、\(f(1)=0\)，并在 \(1\) 的邻域二阶可微且 \(f''(1)\in(0,\infty)\)。

若 \(\tau(t)=t+O(t^{-1})\)，则存在唯一二阶最优轨道
\[
\tau_\star(t)=t+\frac{\sigma^2}{2t}+o(t^{-1})
\]
使
\[
\mathcal D_f(t,\tau_\star(t))
=\frac{f''(1)}{16}\frac{\sigma^4}{t^4}+o(t^{-4}).
\]
并且对任意 \(\tau(t)=t+O(t^{-1})\) 都有
\[
\liminf_{t\to\infty}\frac{t^4}{\sigma^4}\,\mathcal D_f(t,\tau(t))
\ge \frac{f''(1)}{16}.
\]
此外，最优轨道对应的主导似然比形状满足
\[
\frac{t^2}{\sigma^2}\left(\frac{h_t}{g_{t,\tau_\star(t)}}-1\right)(\bar\gamma+ty)
\longrightarrow
-\frac12\,\frac{y^4-6y^2+1}{(1+y^2)^2}
\]
在 \(L^2\!\bigl(\pi^{-1}(1+y^2)^{-1}\dd y\bigr)\) 与一致意义下同时成立。
\end{theorem}
\begin{proof}
记 \(g_t:=g_{t,t}\)，令 \(\tau=t(1+\varepsilon)\) 且 \(\varepsilon=O(t^{-2})\)。
在 \(x=\bar\gamma+ty\) 处有
\[
\frac{g_{t,\tau}}{g_t}
=\frac{(1+\varepsilon)(1+y^2)}{y^2+(1+\varepsilon)^2}
=1+\varepsilon\,S(y)+O(\varepsilon^2),
\qquad
S(y):=\frac{y^2-1}{1+y^2},
\]
其中余项一致于 \(y\in\RR\)。
另一方面，由定理 \ref{thm:xi-cauchy-poisson-second-order-shape-limit-node-rigidity}
\[
\frac{h_t}{g_t}(\bar\gamma+ty)
=1+\frac{\sigma^2}{t^2}U(y)+o(t^{-2}),
\qquad
U(y):=\frac{3y^2-1}{(1+y^2)^2},
\]
一致于 \(y\)。
故
\[
\frac{h_t}{g_{t,\tau}}-1
=\frac{\sigma^2}{t^2}\bigl(U-cS\bigr)+o(t^{-2}),
\qquad
c:=\frac{t^2\varepsilon}{\sigma^2}=O(1),
\]
一致于 \(y\) 且在 \(L^2(G\dd y)\)（\(G(y)=\pi^{-1}(1+y^2)^{-1}\)）中同样成立。

由 \(f\) 在 \(1\) 处二阶展开，
\[
\mathcal D_f(t,\tau)
=\frac{f''(1)}{2}\int_{\RR} g_{t,\tau}
\left(\frac{h_t}{g_{t,\tau}}-1\right)^2\dd x
+o(t^{-4}).
\]
又 \(g_{t,\tau}(\bar\gamma+ty)\,t\,\dd y=G(y)\dd y+o(1)\dd y\)，得
\[
\frac{t^4}{\sigma^4}\mathcal D_f(t,\tau)
=\frac{f''(1)}{2}\,J(c)+o(1),
\qquad
J(c):=\int_{\RR}G(y)\bigl(U(y)-cS(y)\bigr)^2\dd y.
\]
计算
\[
\int GU^2=\frac14,\qquad
\int GUS=\frac14,\qquad
\int GS^2=\frac12,
\]
故
\[
J(c)=\frac14-\frac{c}{2}+\frac{c^2}{2}
=\frac12\left(c-\frac12\right)^2+\frac18.
\]
因此唯一极小点为 \(c_\star=\frac12\)，最小值 \(J(c_\star)=\frac18\)。
于是
\[
\inf_{\tau=t+O(t^{-1})}\liminf_{t\to\infty}\frac{t^4}{\sigma^4}\mathcal D_f(t,\tau)
=\frac{f''(1)}{2}\cdot\frac18
=\frac{f''(1)}{16},
\]
且达到该极限的充要条件为
\(
c(t)\to\frac12
\)，即
\[
\varepsilon(t)=\frac{\sigma^2}{2t^2}+o(t^{-2})
\quad\Longleftrightarrow\quad
\tau_\star(t)=t+\frac{\sigma^2}{2t}+o(t^{-1}).
\]

最后把 \(c_\star=\frac12\) 代回：
\[
U(y)-\frac12S(y)
=-\frac12\,\frac{y^4-6y^2+1}{(1+y^2)^2},
\]
即得最优轨道下的主导似然比形状及其一致/\(L^2(G\dd y)\) 收敛。证毕。
\end{proof}

\endinput
